\chapter{Conclusion and Future Work}

\section{Conclusion}

This project presents a blockchain-based, community-assisted, privacy-preserving reporting framework for local governance. The system is designed to address common problems in centralized reporting platforms, such as lack of transparency, weak data integrity, privacy concerns, and unreliable report quality.

The proposed framework combines a permissioned blockchain, smart contracts, a simulated identity service, a backend relayer, IPFS, AI-assisted moderation, and a web-based DApp. Each component has a clear role in the overall system.

The blockchain and smart contracts provide tamper-evident report records and workflow management. The GovID simulator and ticket-based mechanism support privacy-preserving access control. The relayer improves usability by handling blockchain transactions on behalf of citizens. IPFS supports scalable off-chain storage. The AI oracle service helps filter spam and inappropriate content before permanent recording.

At the current progress stage, several important subsystems have been designed and partially implemented. The AI moderation subsystem has also been deployed as separate oracle containers on a VPS. This shows that the proposed architecture is technically feasible as a prototype.

The main insight from this project is that blockchain alone is not enough for a practical civic reporting system. It must be combined with privacy mechanisms, off-chain storage, moderation, and user-friendly interfaces.

\section{Future Work}

The next major task is to complete full end-to-end integration. The DApp, GovID simulator, backend relayer, AI oracle service, IPFS, smart contracts, and blockchain must work together in one complete report submission flow.

The government-side workflow should also be completed. This includes report review, issue assignment, resolution updates, closure, and citizen feedback.

Image moderation should be added to the AI moderation service. This will allow the system to check uploaded images for harmful, irrelevant, or misleading content.

The current identity system should be improved in future work. A real zero-knowledge proof or self-sovereign identity mechanism can replace the simulated GovID service.

The oracle system can also be improved. At present, decentralization is simulated using separate containers. In future, oracle nodes can be deployed on separate servers or operated by different trusted organizations.

More testing is required after integration. This includes blockchain transaction latency, IPFS upload time, oracle response time, and complete report submission time.

Security testing is also required before any real-world deployment. Smart contracts, APIs, private keys, relayer logic, and deployment infrastructure must be audited carefully.

Finally, the DApp should be improved through usability testing. Citizens should be able to submit and track reports easily without needing blockchain knowledge.

\section{Final Remarks}

The proposed system provides a strong foundation for a transparent and privacy-preserving local governance reporting platform. Although the prototype is still under development, the current progress shows that the architecture is feasible.

With further integration, testing, and improvement, this framework can support more trustworthy citizen participation and better accountability in local governance.